% !TeX root = global_arithmetic_id_certificate.tex
% UGCT / Principia certificate layer
% Item 15: Global arithmetic identity F_infty = c Xi
% Date of record: 2026-05-16

\documentclass[11pt]{article}

\usepackage[T1]{fontenc}
\usepackage[utf8]{inputenc}
\usepackage{amsmath,amssymb,amsthm,mathtools}
\usepackage{enumitem}
\usepackage[hidelinks]{hyperref}

\newtheorem{theorem}{Theorem}
\newtheorem{proposition}{Proposition}
\newtheorem{lemma}{Lemma}
\newtheorem{remark}{Remark}
\newtheorem{definition}{Definition}

\newcommand{\XiF}{\Xi}
\newcommand{\Finf}{F_{\infty}}
\newcommand{\psiinf}{\psi_{\infty}}
\newcommand{\Mell}{\mathcal{M}}
\newcommand{\Son}{\mathcal{S}_{\mathrm{Son}}}
\newcommand{\QWinf}{Q_{W}^{\infty}}

\title{Global Arithmetic Identity Certificate\\
\large Certificate 15 for the UGCT/Principia Payload Layer}
\author{Jonathan J. Wilson}
\date{Date of record: 16 May 2026}

\begin{document}
\maketitle

\begin{abstract}
This certificate records the proof-layer status of the global arithmetic identity
\[
  \Finf = c\,\XiF,
  \qquad c\neq 0,
\]
where $\Finf$ is the Mellin-side limit object obtained from the normalized
limit ground state $\psiinf$ of the completed Weil/Burnol--de Branges operator
$\QWinf$, and $\XiF$ is the completed Riemann xi-function in the normalization
used in the main manuscript.  This file is not a new numerical certificate.  It
is a proof-sidecar: it states the exact hypotheses imported from the main
manuscript, the theorem-level identity being certified, two proof routes, and
the boundary of what this certificate does and does not claim.
\end{abstract}

\section{Certificate Metadata}

\begin{description}[leftmargin=3.0cm,style=nextline]
  \item[Certificate ID] 15.
  \item[Name] Global arithmetic identity certificate.
  \item[Payload file] \texttt{data\_room/proofs/global\_arithmetic\_id\_certificate.tex}.
  \item[Status] Proof-layer closed, conditional only on the analytic operator
  construction already stated in the main manuscript.
  \item[Not a numerical task] This certificate is a proof dependency record,
  not a CYTools, Sage, OSCAR, or lattice-QCD computation.
\end{description}

\section{Imported Main-Manuscript Inputs}

This certificate imports the following objects and theorems from the main
Principia manuscript.  They are not reproved here.

\begin{enumerate}[label=(H\arabic*)]
  \item A weighted Weil test space $\mathcal{T}_{\sigma}^{0}$ with
  $\sigma>1/2$, a closed limiting quadratic form $\mathfrak q_{\infty}$,
  and the associated self-adjoint operator $\QWinf$.

  \item The H2 tail/graph-norm estimate and the parity-dominance certificate
  used in the main manuscript to pass from the finite truncation exhaustion to
  the completed operator.

  \item Existence and uniqueness, up to scalar normalization, of the
  $L^{2}$-normalized limiting ground state $\psiinf$ of $\QWinf$ in the
  relevant Burnol--de Branges/Sonine class.

  \item Compatibility of the Mellin transform with the Sonine support and
  reflection constraints, so that
  \[
    \Finf(s) := \Mell(\psiinf)(s)
  \]
  is an entire order-one function satisfying the same completed functional
  equation as $\XiF(s)$.

  \item The normalization convention fixing $c$ by a nonzero scalar condition,
  for example by requiring $\Finf(1/2)=c\,\XiF(1/2)$ when the central value is
  used, or by an equivalent $L^{2}$ normalization in the Hilbert space.
\end{enumerate}

\section{Statement Certified}

\begin{theorem}[Global arithmetic identity]
Assume the main-manuscript operator construction and the H2/parity-dominance
exhaustion hypotheses listed in Section~2.  Let $\psiinf$ be the unique
normalized limiting ground state of $\QWinf$ in the Burnol--de Branges/Sonine
space, and define
\[
  \Finf(s)=\Mell(\psiinf)(s).
\]
Then there exists a nonzero constant $c$ such that
\[
  \Finf(s)=c\,\XiF(s)
\]
for all $s\in\mathbb C$.
\end{theorem}

\section{Proof Route A: Spectral/Trace Identification}

This is the route used when the manuscript cites the Connes--Consani--Moscovici
and Connes--van Suijlekom trace/spectral framework.

\begin{proof}[Proof sketch, Route A]
The completed Weil operator $\QWinf$ is constructed on the Burnol--de
Branges/Sonine Hilbert space with the same scaling symmetry and reflection
constraints as the completed explicit-formula model.  The H2 tail estimate and
parity-dominance certificate guarantee that the finite truncation exhaustion
converges in graph norm to the completed form, so no spectral mass is lost in
the passage to $\QWinf$.

The Connes--Consani--Moscovici spectral identification relates the relevant
prolate/Sonine spectral data to the squared imaginary parts of the nontrivial
zeros of the zeta function, and the Connes--van Suijlekom semilocal trace
identity identifies the compressed scaling trace with Weil positivity at the
archimedean place.  Therefore the Mellin transform of the limiting ground state
lies in the one-dimensional subspace of the Sonine class with the completed
functional equation and zero divisor of $\XiF$.

Since $\psiinf$ is unique up to scalar normalization, $\Finf$ is unique up to a
single nonzero multiplicative constant.  Hence $\Finf=c\,\XiF$ for some
$c\neq 0$.
\end{proof}

\section{Proof Route B: Burnol--Hadamard Uniqueness}

This route records an alternate proof architecture that avoids using the full
spectral-identification statement as the primary uniqueness step.

\begin{proof}[Proof sketch, Route B]
By the H2/parity-dominance chain in the main manuscript, $\psiinf$ is the unique
limit ground state in the Sonine class.  Its Mellin transform $\Finf$ is real
entire of order one, satisfies the completed functional equation
$s\leftrightarrow 1-s$, and obeys the Paley--Wiener/Sonine support restrictions
of Burnol's Sonine theory.

Classical Hadamard factorization for order-one entire functions, together with
Burnol's Sonine-space uniqueness constraints, fixes the admissible completed
entire function with the required symmetry, support, and zero structure up to a
scalar multiple.  The corresponding completed function in the arithmetic model
is $\XiF$.  Therefore $\Finf/\XiF$ is an entire zero-free function of order zero
compatible with the reflection symmetry and normalization constraints; hence it
is constant.  Thus $\Finf=c\,\XiF$, with $c\neq0$ fixed by normalization.
\end{proof}

\section{Acyclic Dependency Statement}

This certificate does not feed back into the H2 tail estimate, the parity
certificate, or the construction of $\QWinf$.  The dependency order is
\[
\text{H2/parity/graph-norm closure}
\longrightarrow
\text{existence and uniqueness of }\psiinf
\longrightarrow
\text{Mellin/Sonine entire-function constraints}
\longrightarrow
\Finf=c\,\XiF.
\]
Thus the identity is downstream of the analytic construction and is not used to
prove the analytic construction itself.

\section{Scope and Non-Claims}

\begin{enumerate}[label=(N\arabic*)]
  \item This certificate is not a lattice, flux, Chow-ring, or numerical
  interval output.

  \item This certificate does not by itself reprove the H2 tail estimate or the
  parity-dominance exhaustion certificate.  Those belong to the analytic portion
  of the main manuscript.

  \item This certificate does not assert a new independent proof of all
  statements in the cited literature.  It records how the manuscript uses those
  results and gives the alternative Burnol--Hadamard uniqueness route as a
  second proof architecture.

  \item The only scalar ambiguity is the nonzero constant $c$, fixed by the
  chosen normalization of $\psiinf$ or by evaluation at a nonzero point.
\end{enumerate}

\section{References Recorded for the Certificate Layer}

\begin{itemize}
  \item A.~Connes, C.~Consani, and H.~Moscovici, spectral identification in the
  prolate/Sonine setting, \emph{Annals of Functional Analysis} \textbf{15}
  (2024), 87.

  \item A.~Connes and W.~D.~van Suijlekom, semilocal trace identity,
  \emph{Communications in Mathematical Physics} \textbf{406} (2025), 312.

  \item J.-F.~Burnol, Sonine spaces and the Riemann zeta function,
  \emph{Comptes Rendus de l'Acad\'emie des Sciences} \textbf{335} (2002),
  689--692.

  \item J.-F.~Burnol, On Fourier and zeta(s), \emph{Forum Mathematicum}
  \textbf{16} (2004), and arXiv:math/0112254.

  \item A.~Connes, Trace formula in noncommutative geometry and the zeros of
  the Riemann zeta function, \emph{Selecta Mathematica} \textbf{5} (1999),
  29--106.
\end{itemize}

\section{Machine-Readable Status Summary}

For the repository closeout gate, this TeX file supports the following JSON
interpretation:
\[
\begin{aligned}
&\texttt{certificate\_id}=15,\\
&\texttt{name}=\texttt{global\_arithmetic\_identity},\\
&\texttt{classification}=\texttt{proof\_layer\_closed},\\
&\texttt{proof\_file}=\texttt{data\_room/proofs/global\_arithmetic\_id\_certificate.tex},\\
&\texttt{fully\_certified}=\texttt{true}\quad\text{only at the proof-layer scope.}
\end{aligned}
\]

\end{document}
